%% 第三章

\chapter{UAV-DualCog双认知能力测试方法与实验}
\chaptermark{UAV-DualCog双认知能力测试方法与实验}

\section{测试基准设计}

\subsection{总体设计思想}

UAV-DualCog 的核心设计思想是从无人机具身推理需求出发，将多模态大模型能力评测从单一语义问答扩展为“认知维度、观测媒介、证据形式”相结合的系统评测。无人机的观察不是来自被动摄像机，而是由主动运动、飞行高度、航向变化和相机视角共同决定。因此，模型必须同时回答两个问题：一是无人机是否理解自身状态，二是无人机是否理解外部环境。

\begin{figure}[htbp]
    \centering
    \includegraphics[width=\textwidth]{fig_benchmark_overview.pdf}
    \caption{UAV-DualCog测试基准总体结构}
    \label{fig:benchmark-overview}
\end{figure}

如\autoref{fig:benchmark-overview}所示，UAV-DualCog 将自我状态认知和环境情况认知作为两条顶层能力主线，并通过图像任务和视频任务共同测量。图像任务强调离散观察下的空间推理和目标定位，视频任务强调连续观察下的行为理解、可见性判断和时间定位。两类媒介共同服务于同一个问题：模型是否能够在无人机主动观测场景中同时理解“自己”和“外界”，并用可检验的空间或时间证据支撑判断。

\subsection{双认知能力定义}

本文将无人机双认知能力定义为自我状态认知和环境情况认知两部分。对于无人机中心观测 \(o\)，其认知目标可以表示为
\begin{equation}
    \mathcal{C}(o)=\{c_s,c_e\},
    \label{eq:dual-cognition}
\end{equation}
其中，\(c_s\) 表示自我状态认知，\(c_e\) 表示环境情况认知。

自我状态认知关注无人机对自身状态的理解，包括无人机相对于地标的位置、运动导致的未来观察变化、当前执行的飞行行为以及行为发生的时间区间。环境情况认知关注无人机对外部世界的理解，包括地标相对于无人机前向的位置、根据地标位置应采取的动作方向、目标地标在视频中的可见次数和可见时间区间。

这一区分具有实际意义。对于普通视觉问答任务，模型可以只描述图像中有什么目标或目标之间存在什么关系；对于无人机而言，目标外观和空间关系都受到无人机自身位置、朝向和运动的影响。如果模型不能理解自身状态，那么它对外部环境的判断也可能缺乏稳定参照系。因此，双认知能力不是两个互不相关的能力集合，而是无人机具身推理中相互条件化、相互约束的两个方面。

\subsection{六类任务体系}

在双认知框架基础上，本文设计了四类图像任务和两类视频任务。需要说明的是，视频侧的“飞行行为识别与时间定位”在实验统计中进一步拆分为复合行为识别和原子行为识别两个子任务；因此后文视频结果表按复合行为识别、原子行为识别和地标可见性推理三组呈现。六类任务的具体设计如\autoref{tab:task-system}所示。

\begin{table}[htbp]
    \centering
    \caption{UAV-DualCog六类任务体系}
    \label{tab:task-system}
    \begin{tabularx}{\textwidth}{p{3.0cm}p{1.6cm}p{2.0cm}X}
        \toprule
        任务名称 & 模态 & 认知维度 & 任务定义 \\
        \midrule
        地标相对位置推理 & 图像 & 自我状态认知 & 给定地标中心参考视图和当前第一视角观测，判断无人机相对于目标地标的位置，并输出当前图像中的目标框。 \\
        未来观察预测 & 图像 & 自我状态认知 & 给定参考视图和候选未来观察图像，根据动作变化描述选择正确未来观察，并输出对应目标框。 \\
        自我参照位置推理 & 图像 & 环境情况认知 & 在当前第一视角观测中判断目标地标相对于无人机前向的方位，并输出目标框。 \\
        地标驱动动作决策 & 图像 & 环境情况认知 & 根据当前观测中的目标地标位置判断无人机应朝哪个方向移动，并输出目标框。 \\
        飞行行为识别与时间定位 & 视频 & 自我状态认知 & 从第一视角飞行视频中识别无人机正在执行的飞行行为，并给出对应时间区间。 \\
        地标可见次数与区间推理 & 视频 & 环境情况认知 & 判断目标地标在飞行视频中的出现次数，并给出全部可见时间区间。 \\
        \bottomrule
    \end{tabularx}
\end{table}

六类任务在结构上保持统一。图像任务要求输出离散选项和标准化二维目标框；视频任务要求输出行为或可见次数，并给出秒级时间区间。这样设计的原因在于，UAV-DualCog 不希望只评估模型是否能够给出看似正确的自然语言解释，而是要求模型在严格结构化输出下同时完成语义判断和证据定位。

\subsection{输出契约与评测指标}

UAV-DualCog 为图像和视频任务分别设计了严格的结构化输出契约。图像任务统一使用 \texttt{answer\_option\_id} 和 \texttt{bbox\_xyxy\_norm} 字段，视频任务则使用行为答案列表、时间区间或可见次数相关字段。主要评测指标如\autoref{tab:evaluation-protocol}所示。

\begin{table}[htbp]
    \centering
    \caption{任务输出契约与主要评测指标}
    \label{tab:evaluation-protocol}
    \small
    \setlength{\tabcolsep}{4pt}
    \begin{tabularx}{\textwidth}{@{}p{3.0cm}p{2.2cm}p{3.5cm}X@{}}
        \toprule
        任务类型 & 认知维度 & 结构化输出 & 主要指标 \\
        \midrule
        图像任务 & \makecell[l]{自我状态认知\\或环境情况认知} & \makecell[l]{\texttt{answer\_option\_id}\\\texttt{bbox\_xyxy\_norm}} & \makecell[l]{选项准确率\\BBox Acc@0.5\\BBox mIoU} \\
        复合行为识别 & 自我状态认知 & \makecell[l]{\texttt{answers[].option\_id}\\\texttt{answers[].intervals\_sec}} & \makecell[l]{Semantic Accuracy\\Temporal F1@0.5\\Mean tIoU} \\
        原子行为识别 & 自我状态认知 & \makecell[l]{\texttt{answers[].option\_id}\\\texttt{answers[].intervals\_sec}} & \makecell[l]{Semantic Accuracy\\Temporal F1@0.5\\Mean tIoU} \\
        地标可见性推理 & 环境情况认知 & \makecell[l]{\texttt{visible\_count}\\\texttt{visible\_intervals\_sec}} & \makecell[l]{Count Accuracy\\Segment F1@0.5\\Mean tIoU} \\
        \bottomrule
    \end{tabularx}
\end{table}

从评测协议可以看出，本文将目标框和时间区间作为正式指标，而不是附属可视化信息。这使得实验可以进一步分析模型是否存在“判断正确但证据错误”的现象。例如，图像任务中模型可能选对方向选项，却无法定位地标；视频任务中模型可能判断出目标出现次数，却无法给出完整可见区间。UAV-DualCog 通过将这些证据纳入指标，能够更真实地衡量无人机具身推理的可靠性。

\subsection{提示词协议设计}

为了保证不同模型在统一接口下可比较，本文对六类任务设计了固定提示词协议。每个任务由系统提示词、用户提示模板和 JSON 输出模式三部分组成。系统提示词规定无人机代理身份、输入类型、输出字段和禁止事项；用户提示词注入当前任务中的地标描述、参考图像、当前观测、候选选项或视频语境；JSON 输出模式则保证解析器可以稳定提取字段并计算指标。正文给出协议摘要，完整协议结构见附录\autoref{app:prompt-protocol}。

\begin{table}[htbp]
    \centering
    \caption{六类任务提示词设计摘要}
    \label{tab:prompt-summary}
    \begin{tabularx}{\textwidth}{p{3.1cm}X p{3.5cm}}
        \toprule
        任务名称 & 核心约束 & 输出契约 \\
        \midrule
        地标相对位置推理 & 指定无人机代理身份，输入参考视图和当前视图，必须同时给出选项与目标框。 & \texttt{\{answer, bbox\}} \\
        未来观察预测 & 从多个候选未来视图中选择正确观察，并给出选中图像中的目标框。 & \texttt{\{answer, bbox\}} \\
        自我参照位置推理 & 判断目标地标相对无人机前向方向，并在当前观测中定位目标。 & \texttt{\{answer, bbox\}} \\
        地标驱动动作决策 & 根据地标位置选择接近目标的飞行动作，并输出作为证据的目标框。 & \texttt{\{answer, bbox\}} \\
        飞行行为识别与时间定位 & 识别一个或多个飞行行为选项，并给出各自对应时间区间。 & \texttt{\{answers\}} \\
        地标可见次数与区间推理 & 输出目标地标可见次数和全部可见时间区间。 & \texttt{\{count, intervals\}} \\
        \bottomrule
    \end{tabularx}
\end{table}

这种协议层设计将答案字段和证据字段绑定在一起，避免模型通过宽泛自然语言描述绕开评测解析。对于后续能力优化而言，结构化输出也为监督微调和可验证奖励提供了基础。

\subsection{测试基准规模与统计}

UAV-DualCog 的完整资产池覆盖 18 个 AerialVLN 场景，总面积约为 5040232.00 m\(^2\)。经过自动构建和人工复核后，从 2626 个候选地标中保留 746 个有效地标，覆盖 8 个粗粒度类别和 166 个细粒度子类。当前发布的测试基准从中选择 12 个场景和 512 个有效地标，构建 4096 个图像任务样本和 2048 个视频任务样本。

图像子集包含四类任务，每类 1024 个样本，共使用 3826 张唯一源图像和 12288 次图像引用。视频子集包含两类任务，每类 1024 个样本，基于 1024 条飞行轨迹生成，总轨迹长度为 388661.79 m，总视频时长为 19633.30 s，约合 5 小时 27 分 13 秒。

\begin{figure}[htbp]
    \centering
    \includegraphics[width=\textwidth]{fig_benchmark_stat.pdf}
    \caption{UAV-DualCog测试基准统计信息}
    \label{fig:benchmark-stat}
\end{figure}

如\autoref{fig:benchmark-stat}所示，测试基准在任务数量、认知维度和媒介类型上保持均衡。每类任务均为 1024 个样本，可以避免某一任务因样本数量过大而主导总体评测结果。场景方面覆盖城市、乡镇、工业区、水滨区域、夜间场景等多种低空无人机常用环境，为模型能力分析提供了较丰富的场景基础。

\subsection{测试基准设计小结}

本节介绍了 UAV-DualCog 的总体设计、双认知定义、六类任务体系、输出契约、提示词协议和测试基准规模。该基准的关键特点在于将自我状态认知与环境情况认知同时纳入评测，并把空间目标框和时间区间作为正式证据指标。下文将进一步介绍测试基准数据构建工具链的实现过程。

\FloatBarrier
\section{测试基准数据构建工具链实现}

\subsection{工具链总体流程}

UAV-DualCog 的数据构建并非一次性人工整理，而是由一套多阶段工具链完成。工具链的总体目标是从仿真场景出发，经过场景几何建模、地标资产构建、视频任务生成和图像任务生成，最终导出可评测、可复核、可追踪的结构化测试基准。

\begin{figure}[htbp]
    \centering
    \includegraphics[width=\textwidth]{fig_benchmark_construction_process.pdf}
    \caption{UAV-DualCog四阶段构建流程}
    \label{fig:construction-process}
\end{figure}

如\autoref{fig:construction-process}所示，工具链分为四个阶段：阶段1完成场景扫描与点云融合，阶段2完成地标挖掘、人工复核与语义标注，阶段3完成行为驱动的视频任务构建，阶段4完成证据感知的图像问答任务构建。各阶段与仓库真实实现入口、关键伪代码和代表性逻辑的对应关系见\autoref{tab:app-toolchain}。正文重点说明其数据流、关键约束和任务产物，避免将工程实现直接写成文件清单。

\subsection{阶段1：场景扫描与点云融合}

阶段1的目标是建立稳定的场景级几何基础。该阶段基于 AerialVLN 模拟器场景和配置文件，在给定边界、采样高度和姿态覆盖范围内进行多位姿扫描，同时采集 RGB 图像、语义分割、激光雷达和无人机位姿信息。随后，系统将这些原始数据统一变换到同一场景坐标系中，融合为带有语义和实例属性的场景级点云。

这一阶段的关键作用在于将原始模拟器场景转换为可复用的几何基底。完成阶段1后，每个场景具备稳定的地图边界、全局坐标系、语义点云、扫描预览图和面积统计信息。后续阶段不需要重复依赖局部图像检测，而可以直接在场景级几何表示上开展地标候选聚合、可见性判断和视角采样。

从工程实现上，阶段1围绕场景配置解析、模拟器桥接采集和场景级产物落盘三个环节组织。系统统一读取场景配置，调用模拟器桥接模块完成多位姿采样，并将点云、日志和缓存结果写入标准化目录结构中。具体实现入口与伪代码见\autoref{tab:app-toolchain}。

\subsection{阶段2：地标挖掘、复核与语义标注}

阶段2基于阶段1输出的语义与实例点云构建地标资产。系统首先按照实例信息聚合候选对象，形成初始地标候选池；随后根据候选对象三维包围盒和实例中心，计算八个标准方向的拍摄位姿，包括前、右前、右、右后、后、左后、左、左前等视角，并生成对应 RGB 视图和点云视图。

自动候选生成之后，项目引入人工复核机制。复核内容包括候选是否保留为有效地标、主视图是否合理、方向标注是否正确、某些视角是否需要剔除、地标框是否有效以及该地标是否具备后续任务构建价值。通过人工复核，阶段2输出形成有效地标清单，作为图像任务和视频任务的共同资产基础。

在语义标注方面，项目使用受约束的自动标注流程，为每个地标补充粗粒度类别、细粒度子类和简短可区分描述。与自由生成式描述不同，该流程通过系统提示词、用户提示模板和固定 JSON 输出格式约束模型生成，使地标语义信息既可读又可解析。该设计为后续任务提示词、样本复核和实验记录提供了统一语义基础。

阶段2的主要实现入口见\autoref{tab:app-toolchain}。该阶段支持实例收集、复核界面、语义标注和有效实例导出等模式。通过复核界面，标注者可以同时查看实例级信息、RGB 视图、点云视图、方向标签和语义描述，从而提高地标质量控制效率。

\subsection{阶段3：行为驱动的视频任务构建}

阶段3负责视频任务生成。该阶段基于阶段2得到的有效地标资产和预设行为库，为地标绑定飞行任务，并导出轨迹、视频、行为区间和可见性监督信息。

该阶段的一个重要设计是两层飞行行为体系。底层为原子行为，例如渐进接近、环绕飞行、螺旋环绕、升空观察和面状扫描等；上层为复合行为，例如环绕巡检、螺旋巡检、三角巡检和面状测绘巡检等。复合行为由多个原子行为组合而成，更接近真实巡检和测绘任务中的飞行模式。

阶段3实现包括任务模板参数确定、轨迹生成、碰撞检查、安全距离约束、半径和高度修正、并行调度、关键帧与时间区间监督导出等步骤。最终，该阶段保留高分辨率原始帧和任务级监督文件，同时导出面向测试基准的 1080P H.264 视频和结构化任务清单，使视频任务既适合模型调用，又保留足够的时间定位监督。

从实现结构看，阶段3由轨迹生成与任务组织两个子流程共同承担，具体入口和关键伪代码见\autoref{tab:app-toolchain}。前者负责生成飞行轨迹、视频和过程记录，后者负责组织任务、实验和指标聚合。核心视频任务分别覆盖自我状态认知下的飞行行为识别与时间定位，以及环境情况认知下的地标可见次数与区间推理。

\subsection{阶段4：证据感知图像问答任务构建}

阶段4负责图像任务生成。该阶段的核心思想不是随机选取截图组成问答，而是围绕地标中心参考视图和任务相关第一视角观测图进行结构化样本生成。自我状态认知和环境情况认知两类图像任务共享部分地标视图资产，但分别围绕“自身位置与未来观察”和“目标方向与动作决策”展开问题设计。

图像任务生成中，系统同时考虑语义答案和空间证据。对于需要定位的任务，样本除问题、选项和答案外，还包含标准化目标框。对于环境情况认知任务，系统会调整观测位姿和相机朝向，直到目标完整进入画面且目标框有效为止。最终输出不仅包括任务清单，也包括重新渲染所需辅助文件，以便在不改变样本语义结构的前提下刷新媒体资源。

阶段4的主要实现入口见\autoref{tab:app-toolchain}。该阶段支持图像 QA 生成、任务渲染、模型调用、结构化输出解析和指标汇总。四类图像任务分别覆盖自我状态认知与环境情况认知下的两个细分方向，并统一使用选项答案和目标框作为输出。

\subsection{统一管线与配套材料整理}

为了降低多阶段构建带来的复杂度，项目实现了统一管线编排。仓库中的 \texttt{task\_pipeline.py} 支持跨场景的数据筛选、任务生成、渲染、实验和分析流程；\texttt{pipeline\_common.py} 则统一桥接配置、场景目录解析、阶段目录解析、完整性校验和阶段日志。典型场景目录包含 \texttt{pcd\_map}、\texttt{landmarks\_raw}、\texttt{landmarks\_review}、\texttt{landmarks}、\texttt{stage3\_tasks} 和 \texttt{logs} 等子目录。

配套材料方面，项目整理了 README、数据集说明、使用说明、实验榜单、分析结果和配置模板。这些材料分别对应基准介绍、数据使用、模型评测和结果分析等环节，用于支撑测试基准的复核、复现实验和后续扩展。

\subsection{数据构建小结}

本节介绍了 UAV-DualCog 测试基准的四阶段构建工具链。阶段1提供场景几何基础，阶段2提供复核后的地标资产，阶段3生成视频任务和时间监督，阶段4生成图像任务和空间证据。通过共享配置、统一路径和结构化输出，该工具链既支撑了当前测试基准的稳定构建，也已经实际支撑 UAV-DualCog-Train 训练集的数据导出。工具链实现入口与关键伪代码见\autoref{tab:app-toolchain}。

\FloatBarrier
\section{能力评测与结果分析}

\subsection{实验设置}

在测试基准构建完成后，本文进一步完成统一评测协议设计和多模型实验。UAV-DualCog 不采用简单的答案字符串匹配方式，而是要求所有任务输出严格 JSON，并分别计算语义指标和证据指标。也就是说，实验不仅关注模型是否回答正确，还关注模型是否能够用正确的空间或时间证据支撑回答。

对于图像任务，模型需要预测离散选项和标准化目标框。本文报告选项准确率、BBox Acc@0.5 和 BBox mIoU。对于视频任务，模型需要预测行为或可见性结果，并给出时间区间。本文报告复合行为准确率、原子行为准确率、可见次数准确率、Temporal F1@0.5 和 Mean tIoU 等指标。图像样本以压缩到 \(640\times480\) 的 JPEG 输入为主，视频样本以压缩 MP4 视频和必要参考图像作为输入。

为减少不同模型家族在显式思维设置上的差异，实验统一采用 Instant 或等价的快速推理模式，关闭额外显式推理开关，使实验更接近无人机实时感知场景。实验模型覆盖闭源模型、开源模型和空间推理微调模型三类。模型来源如\autoref{tab:model-source}所示。

\begin{table}[htbp]
    \centering
    \caption{实验模型类型与来源}
    \label{tab:model-source}
    \begin{tabularx}{\textwidth}{p{2.2cm}p{2.7cm}X p{2.6cm}}
        \toprule
        模型类型 & 模型家族 & 具体型号 & 来源 \\
        \midrule
        闭源模型 & Claude、GPT、Gemini、Grok、Qwen商业版、Mimo & Claude Sonnet 4.6、GPT 5.3 Chat、Gemini 3 Flash、Gemini 3.1 Flash Lite、Grok 4.1 Fast、Qwen 3.5/3.6 系列、Mimo v2 Omni & OpenRouter、Bailian、Xiaomi Mimo \\
        开源模型 & GLM、Kimi、Qwen、InternVL、Intern-S1 & GLM 4.6V、Kimi K2.5、Qwen 3.5 系列、InternVL 3.5 系列、Intern-S1 Pro & Zhipu AI、Siliconflow、本地部署、SH-AILab \\
        空间推理模型 & SpaceR、SpaceThinker、SpaceOm、SenseNova、ViLaSR、VST & SpaceR、SpaceThinker、SpaceOm、SenseNova-SI-1.2、ViLaSR、VST-7B-RL、VST-7B-SFT & 本地部署 \\
        \bottomrule
    \end{tabularx}
\end{table}

为便于结果解读，本文同时保留代表性模型版本对应的公开技术来源。部分闭源模型名称以实验调用平台在实验日期返回的模型标识为准；对于尚无正式公开论文的模型，本文仅将其作为实验接口来源记录，不将其技术报告作为本文方法设计依据。闭源模型覆盖 Claude Sonnet 4.6、GPT 5.3、Gemini 3 Flash、Gemini 3.1 Flash Lite、Grok 4.1 Fast、Qwen 3.5/3.6 商业版与 Mimo v2 Omni\cite{anthropic2026claude46,openai2026gpt53,deepmind2025gemini3flash,deepmind2026gemini31flashlite,xai2025grok41,qwen3.5,qwen36plus,mimo_v2_omni}；开源模型覆盖 GLM 4.6V、Kimi K2.5、InternVL 3.5 与 Intern-S1 Pro\cite{vteam2025glm45vglm41vthinkingversatilemultimodal,team2026kimi,wang2025internvl3_5,zou2026interns1proscientificmultimodalfoundation}；空间推理模型覆盖 SpaceR、SpaceThinker、SpaceOm、SenseNova-SI-1.2、ViLaSR 与 VST 系列\cite{ouyang2025spacer,SpaceThinker2025,SpaceOm2025,sensenova-si,wu2025reinforcing,yang2025visual}。

本部分按照两个核心问题展示实验结果和分析：第一，图像与视频两种媒介下的双认知能力是否能够相互迁移；第二，自我状态认知与环境情况认知两条能力主线是否已经发展均衡。为保证实验结果、对应分析和交叉结论在正文中完整闭环呈现，本部分直接给出各类核心实验大表、相关性统计表和关键分析图；六类任务提示词协议见\autoref{app:prompt-protocol}。

\definecolor{bestblue}{HTML}{60CBF3}
\definecolor{secondblue}{HTML}{CBEEFB}
\newcommand{\resultgroup}[2]{\midrule\multicolumn{#1}{c}{#2} \\ \midrule}

\subsection{跨媒介迁移分析：双认知能力在图像与视频之间存在真实但不充分的迁移}

本文将跨媒介迁移定义为图像总体指标与视频总体指标之间存在正相关，即图像侧较强的模型在视频侧通常也不会完全失效；同时，将迁移不充分定义为同一模型在两种媒介上的得分差异较大，且散点分布没有沿对角线稳定聚集。基于这一定义，后文从图像任务、视频任务和跨媒介综合统计三个层面展开分析。

UAV-DualCog 的一个重要设计是，同一套双认知能力并不只在单一媒介下测试，而是在图像和视频两种证据通道中分别测量。图像任务强调离散观察下的空间推理和目标定位，视频任务强调连续观察下的行为理解、可见性判断和时间定位。因而，跨媒介分析的重点不是简单比较图像难还是视频难，而是考察模型能否在不同媒介下保持一致、稳定、可证据验证的双认知能力。

\subsubsection{图像任务结果与分析}

图像任务包括地标相对位置推理、未来观察预测、自我参照位置推理和地标驱动动作决策四类任务，覆盖自我状态认知和环境情况认知两条能力主线。由于图像任务同时呈现选项准确率和空间定位指标，因此能够直接观察模型是否只是选项判断正确，还是能够进一步给出可信的空间证据。\autoref{tab:image-results}汇总了全部已评测模型的图像总体准确率、自我状态认知准确率、环境情况认知准确率和代表性空间证据均值。

\begin{table}[htbp]
    \centering
    \caption{图像任务完整结果表}
    \label{tab:image-results}
    \scriptsize
    \renewcommand{\arraystretch}{0.9}
    \setlength{\tabcolsep}{4pt}
    \resizebox{\textwidth}{!}{%
    \begin{tabular}{lrrrrrrrrrrrrr}
        \toprule
        模型 & 全局准确率 & \multicolumn{6}{c}{自我状态认知} & \multicolumn{6}{c}{环境情况认知} \\
         &  & \multicolumn{3}{c}{地标相对位置推理} & \multicolumn{3}{c}{未来观察预测} & \multicolumn{3}{c}{自我参照位置推理} & \multicolumn{3}{c}{地标驱动动作决策} \\
         &  & Acc & BBox Acc@0.5 & mIoU & Acc & BBox Acc@0.5 & mIoU & Acc & BBox Acc@0.5 & mIoU & Acc & BBox Acc@0.5 & mIoU \\
        \midrule
        \resultgroup{14}{闭源模型}
        Claude Sonnet 4.6 & 42.7\% & 37.7\% & 4.9\% & 18.2\% & 23.6\% & 3.4\% & 14.2\% & 48.5\% & 9.6\% & 20.2\% & 61.0\% & 8.6\% & 19.9\% \\
        GPT 5.3 Chat & \cellcolor{secondblue}{47.8\%} & 35.2\% & 11.9\% & 23.5\% & 37.3\% & 12.4\% & 25.2\% & 56.5\% & 14.9\% & 22.5\% & 62.3\% & 15.4\% & 22.5\% \\
        Gemini 3 Flash & \cellcolor{bestblue}{50.2\%} & 47.6\% & 0.7\% & 1.2\% & 45.9\% & 0.7\% & 0.9\% & 56.2\% & 0.1\% & 0.9\% & 51.3\% & 0.5\% & 1.2\% \\
        Gemini 3.1 Flash Lite & 34.5\% & 39.1\% & 0.0\% & 0.0\% & 28.6\% & 2.5\% & 3.5\% & 39.9\% & 0.4\% & 0.8\% & 30.4\% & 0.4\% & 0.6\% \\
        Grok 4.1 Fast & 27.4\% & 21.1\% & 3.6\% & 16.4\% & 22.9\% & 2.2\% & 17.5\% & 33.0\% & 1.9\% & 8.2\% & 32.6\% & 1.6\% & 7.5\% \\
        Qwen 3.6-Plus & 39.3\% & 32.8\% & 16.6\% & 27.1\% & 33.5\% & 14.0\% & 26.7\% & 48.4\% & 7.4\% & 19.0\% & 42.5\% & 6.0\% & 19.9\% \\
        Qwen 3.5-Plus & 38.5\% & 28.6\% & 21.4\% & 27.1\% & 29.0\% & 13.4\% & 23.5\% & 47.6\% & 9.7\% & 19.3\% & 48.9\% & 8.3\% & 20.0\% \\
        Qwen 3.5-Flash & 40.9\% & 27.8\% & 21.6\% & 29.8\% & 28.9\% & 15.9\% & 23.3\% & 52.4\% & 8.8\% & 19.4\% & 54.5\% & 8.9\% & 20.1\% \\
        Mimo v2 Omni & 31.2\% & 28.7\% & 38.0\% & 34.3\% & 21.3\% & 29.8\% & 29.7\% & 36.5\% & 8.8\% & 17.4\% & 38.2\% & 9.2\% & 17.3\% \\
        \resultgroup{14}{开源模型}
        GLM 4.6V & 32.8\% & 30.8\% & 9.4\% & 8.7\% & 27.1\% & 8.6\% & 8.7\% & 29.2\% & 3.3\% & 6.7\% & 44.2\% & 2.5\% & 5.0\% \\
        Kimi K2.5 & 34.1\% & 34.3\% & 27.1\% & 30.4\% & 32.7\% & 15.0\% & 21.7\% & 39.4\% & 5.0\% & 15.6\% & 30.1\% & 4.7\% & 14.4\% \\
        Qwen 3.5-397B-A17B & 39.4\% & 27.4\% & 3.3\% & 12.6\% & 35.8\% & 10.0\% & 21.4\% & 47.9\% & 2.8\% & 13.0\% & 46.2\% & 4.4\% & 13.9\% \\
        Qwen 3.5-122B-A10B & \cellcolor{secondblue}{42.3\%} & 32.6\% & 29.1\% & 31.6\% & 37.8\% & 20.6\% & 28.4\% & 49.7\% & 6.1\% & 19.4\% & 49.0\% & 6.8\% & 18.8\% \\
        Qwen 3.5-35B-A3B & 41.9\% & 29.5\% & 27.6\% & 30.5\% & 27.4\% & 0.5\% & 0.8\% & 53.2\% & 6.1\% & 14.5\% & 57.5\% & 7.2\% & 18.5\% \\
        Qwen 3.5-27B & \cellcolor{bestblue}{44.3\%} & 31.6\% & 42.8\% & 39.9\% & 37.1\% & 17.1\% & 25.8\% & 57.8\% & 7.3\% & 20.4\% & 50.6\% & 7.5\% & 21.0\% \\
        Qwen 3.5-9B & 40.9\% & 29.0\% & 15.9\% & 25.8\% & 29.0\% & 17.8\% & 29.8\% & 53.2\% & 10.3\% & 20.9\% & 52.2\% & 11.5\% & 22.1\% \\
        Qwen 3.5-4B & 39.0\% & 30.9\% & 12.8\% & 23.1\% & 30.2\% & 11.2\% & 21.7\% & 47.5\% & 6.7\% & 17.7\% & 47.5\% & 6.3\% & 18.0\% \\
        Intern S1-Pro & 28.4\% & 29.4\% & 17.4\% & 26.8\% & 26.2\% & 9.6\% & 21.3\% & 27.8\% & 2.7\% & 13.7\% & 30.3\% & 3.2\% & 14.0\% \\
        InternVL 3.5-241B-A28B & 37.7\% & 32.9\% & 28.4\% & 31.2\% & 26.1\% & 19.5\% & 25.3\% & 50.1\% & 5.4\% & 15.8\% & 41.8\% & 5.1\% & 15.0\% \\
        InternVL 3.5-30B-A3B & 32.0\% & 27.3\% & 4.9\% & 19.9\% & 24.0\% & 7.0\% & 15.5\% & 35.4\% & 0.6\% & 7.6\% & 41.3\% & 1.3\% & 7.3\% \\
        InternVL 3.5-14B & 31.1\% & 22.8\% & 7.6\% & 22.2\% & 31.9\% & 6.6\% & 22.0\% & 28.5\% & 4.7\% & 12.9\% & 41.3\% & 5.0\% & 11.9\% \\
        InternVL 3.5-8B & 27.1\% & 28.5\% & 19.7\% & 34.2\% & 24.5\% & 8.1\% & 28.7\% & 24.7\% & 4.2\% & 12.3\% & 30.8\% & 2.5\% & 11.4\% \\
        InternVL 3.5-4B & 27.8\% & 25.6\% & 3.1\% & 19.8\% & 26.4\% & 10.8\% & 23.2\% & 22.4\% & 1.0\% & 8.0\% & 36.7\% & 1.0\% & 7.5\% \\
        \resultgroup{14}{空间推理模型}
        SpaceR & 34.0\% & 20.6\% & 2.4\% & 6.0\% & 24.0\% & 1.2\% & 1.9\% & 42.4\% & 1.8\% & 5.4\% & 48.8\% & 2.3\% & 6.8\% \\
        SpaceThinker & 28.8\% & 22.6\% & 1.2\% & 3.0\% & 25.4\% & 1.9\% & 5.7\% & 32.8\% & 0.6\% & 5.1\% & 34.3\% & 0.7\% & 4.7\% \\
        SpaceOm & 29.3\% & 24.7\% & 2.8\% & 6.7\% & 24.5\% & 3.6\% & 10.9\% & 35.2\% & 2.4\% & 8.6\% & 32.8\% & 1.8\% & 6.5\% \\
        SenseNova-SI-1.2 & 19.6\% & 17.5\% & 10.7\% & 10.6\% & 0.1\% & 0.0\% & 0.1\% & 15.3\% & 3.2\% & 5.8\% & 45.4\% & 6.0\% & 8.5\% \\
        ViLaSR & 33.4\% & 18.8\% & 0.0\% & 0.0\% & 22.1\% & 0.0\% & 0.0\% & 43.1\% & 0.0\% & 0.0\% & 49.7\% & 0.0\% & 0.0\% \\
        VST-7B-RL & \cellcolor{bestblue}{36.1\%} & 29.2\% & 0.0\% & 0.0\% & 24.1\% & 0.0\% & 0.0\% & 47.2\% & 0.0\% & 0.1\% & 44.0\% & 0.0\% & 0.0\% \\
        VST-7B-SFT & \cellcolor{secondblue}{36.1\%} & 28.1\% & 0.0\% & 0.0\% & 23.0\% & 0.0\% & 0.0\% & 46.5\% & 0.0\% & 0.2\% & 46.6\% & 0.0\% & 0.2\% \\
        \bottomrule
      \end{tabular}%
    }
\end{table}

从\autoref{tab:image-results}可以看出，图像任务结果中 Gemini 3 Flash、GPT 5.3 Chat 和多个 Qwen 3.5 系列模型表现突出，但它们的优势并不完全相同。Gemini 3 Flash 的图像总体准确率最高，达到 50.2\%，说明其选项层面的语义判断能力较强；但该模型平均 BBox Acc@0.5 仅约 0.5\%，说明强语义判断并没有自然转化为精确空间证据。GPT 5.3 Chat 的图像总体准确率为 47.8\%，平均 BBox Acc@0.5 达到 13.7\%，空间证据质量相对更均衡。Qwen 3.5-27B 的图像总体准确率为 44.3\%，其中地标相对位置推理 BBox Acc@0.5 达到 42.77\%，体现出较好的局部空间定位潜力。相比之下，VST-7B-RL、VST-7B-SFT、ViLaSR 等空间推理或强化模型虽然在部分选项准确率上并非完全失效，但 BBox 指标大量接近 0，说明一般空间推理微调并不能直接解决 UAV-DualCog 所要求的“选项判断+显式证据”联合输出问题。

图像分支的核心问题是：模型即使能够在自我状态认知或环境情况认知任务中选出正确答案，是否也能同时恢复出支撑该判断的空间证据。选取各模型家族的代表性模型进行统计，图像任务中的平均语义成绩达到 35.98\%，而 BBox Acc@0.5 的平均值仅为 7.02\%。这说明当前模型在图像分支上已经能够形成一定程度的双认知语义判断，但这种判断远未能稳定落实为可验证的目标定位结果。

\begin{figure}[htbp]
    \centering
    \includegraphics[width=0.85\textwidth]{fig_analysis_image_semantic_vs_spatial.png}
    \caption{图像任务语义准确率与空间定位准确率对比}
    \label{fig:image-semantic-spatial}
\end{figure}

如\autoref{fig:image-semantic-spatial}所示，图像分支中“语义正确但证据不足”是一个高度普遍的现象。GPT 5.3 Chat、Claude Sonnet 4.6、Gemini 3 Flash 等模型在选项准确率上都能保持较强竞争力，但在目标框对齐上仍然存在较大损失。相比之下，Mimo v2 Omni 在图像总榜中的总体准确率并非最高，却表现出相对更强的空间证据恢复能力，说明不同模型家族在“判断正确”和“证据对齐正确”之间存在明显不同的偏向。Qwen 3.5 系列则整体表现出较好的稳定性，在多项图像任务上同时兼顾语义判断和一定程度的空间定位能力。

从任务级相关性看，图像分支内部并非所有任务都暴露出同样程度的证据瓶颈。\autoref{tab:image-correlation}显示，地标驱动动作决策任务的语义-空间相关性最高，Pearson \(r=0.44\)；而地标相对位置推理任务的相关性仅为 0.02，几乎不表现出稳定耦合关系。

\begin{table}[htbp]
    \centering
    \caption{图像分支任务级相关性分析表}
    \label{tab:image-correlation}
    \begin{tabularx}{\textwidth}{Xrrr}
        \toprule
        任务 & Pearson \(r\) & 语义均值 & 空间证据均值 \\
        \midrule
        地标相对位置推理 & 0.02 & 30.35 & 10.69 \\
        未来观察预测 & 0.15 & 26.29 & 7.81 \\
        自我参照位置推理 & 0.28 & 41.97 & 4.73 \\
        地标驱动动作决策 & 0.44 & 45.33 & 4.85 \\
        \bottomrule
    \end{tabularx}
\end{table}

\autoref{tab:image-correlation}说明，四类图像任务都存在语义与空间证据不完全同步的问题，其中环境情况认知任务的相关性略高，而需要模型显式处理无人机相对地标位置关系的任务相关性最低。尤其是地标相对位置推理的语义-空间相关性仅为 0.02，说明自我状态认知推理中的空间证据恢复更具挑战性。这意味着，在无人机场景中，“选项判断正确”很可能只是利用了语言先验、局部外观或候选项偏差，并不必然代表模型已经完成了可验证的空间推理，也不能直接视为有确切证据支撑的可信结论。

\FloatBarrier
\subsubsection{视频任务结果与分析}

视频任务包括复合行为识别与时间定位、原子行为识别与时间定位、地标可见次数与区间推理，分别考察模型对无人机自身飞行行为和外部目标可见过程的理解。与图像任务中的空间目标框不同，视频任务中的证据形式是时间区间，因此能够检验模型是否能够把语义判断落到具体时间证据上。\autoref{tab:video-results}汇总了全部已评测模型的视频结果。

\begin{table}[htbp]
    \centering
    \caption{视频任务完整结果表}
    \label{tab:video-results}
    \scriptsize
    \renewcommand{\arraystretch}{0.9}
    \setlength{\tabcolsep}{4pt}
    \resizebox{\textwidth}{!}{%
    \begin{tabular}{lrrrrrrrrrr}
        \toprule
        模型 & 全局准确率 & \multicolumn{3}{c}{复合行为识别} & \multicolumn{3}{c}{原子行为识别} & \multicolumn{3}{c}{地标可见性推理} \\
         &  & Acc & F1@0.5 & Mean tIoU & Acc & F1@0.5 & Mean tIoU & Acc & F1@0.5 & Mean tIoU \\
        \midrule
        \resultgroup{11}{闭源模型}
        Gemini 3 Flash & \cellcolor{bestblue}{46.5\%} & 49.3\% & 51.9\% & 54.3\% & 35.0\% & 52.6\% & 46.6\% & 55.2\% & 44.9\% & 44.3\% \\
        Gemini 3.1 Flash Lite & 31.9\% & 22.3\% & 29.5\% & 32.2\% & 19.7\% & 40.0\% & 40.6\% & 53.7\% & 43.1\% & 41.6\% \\
        Mimo v2 Omni & \cellcolor{secondblue}{38.8\%} & 27.8\% & 33.1\% & 35.8\% & 29.1\% & 53.6\% & 46.9\% & 59.4\% & 52.5\% & 53.7\% \\
        Qwen 3.5-Flash & 36.8\% & 32.6\% & 35.8\% & 40.1\% & 25.6\% & 51.7\% & 42.4\% & 52.3\% & 57.8\% & 51.3\% \\
        Qwen 3.5-Plus & 28.6\% & 2.4\% & 9.0\% & 9.2\% & 30.0\% & 46.6\% & 47.7\% & 53.2\% & 53.9\% & 46.8\% \\
        Qwen 3.6-Plus & 29.6\% & 9.8\% & 16.5\% & 16.9\% & 28.3\% & 45.9\% & 47.3\% & 50.8\% & 51.5\% & 46.1\% \\
        \resultgroup{11}{开源模型}
        GLM 4.6V & 32.4\% & 32.4\% & 1.3\% & 6.7\% & 33.4\% & 21.6\% & 22.6\% & 31.3\% & 16.7\% & 33.8\% \\
        Kimi K2.5 & 30.9\% & 11.9\% & 19.1\% & 19.9\% & 30.2\% & 54.6\% & 49.2\% & 50.5\% & 49.7\% & 45.9\% \\
        Qwen 3.5-397B-A17B & 23.4\% & 5.2\% & 1.7\% & 2.2\% & 24.6\% & 29.9\% & 30.0\% & 40.4\% & 20.4\% & 31.8\% \\
        Qwen 3.5-122B-A10B & 30.6\% & 21.2\% & 28.4\% & 31.9\% & 22.1\% & 39.0\% & 34.6\% & 48.4\% & 46.4\% & 49.7\% \\
        Qwen 3.5-35B-A3B & 34.0\% & 21.6\% & 21.5\% & 23.1\% & 26.2\% & 54.4\% & 50.3\% & 54.1\% & 49.7\% & 48.6\% \\
        Qwen 3.5-27B & 31.3\% & 14.6\% & 21.9\% & 23.3\% & 28.0\% & 54.5\% & 50.7\% & 51.4\% & 51.7\% & 48.7\% \\
        Qwen 3.5-9B & 29.2\% & 9.9\% & 10.1\% & 11.0\% & 27.2\% & 55.9\% & 51.3\% & 50.7\% & 33.5\% & 36.8\% \\
        Qwen 3.5-4B & 16.5\% & 6.4\% & 1.3\% & 1.4\% & 9.4\% & 3.0\% & 2.8\% & 33.9\% & 44.3\% & 52.0\% \\
        InternVL 3.5-38B & \cellcolor{bestblue}{37.8\%} & 15.2\% & 11.1\% & 12.8\% & 34.2\% & 33.0\% & 33.1\% & 64.0\% & 20.0\% & 27.5\% \\
        InternVL 3.5-30B-A3B & 28.0\% & 15.3\% & 5.8\% & 8.6\% & 24.8\% & 24.2\% & 24.6\% & 43.9\% & 7.9\% & 12.6\% \\
        InternVL 3.5-14B & \cellcolor{secondblue}{36.6\%} & 33.1\% & 3.9\% & 11.3\% & 22.4\% & 15.8\% & 17.4\% & 54.5\% & 13.9\% & 22.7\% \\
        InternVL 3.5-8B & 34.1\% & 27.4\% & 8.5\% & 14.1\% & 23.0\% & 25.0\% & 26.1\% & 51.9\% & 16.4\% & 21.3\% \\
        InternVL 3.5-4B & 28.9\% & 31.0\% & 3.5\% & 7.4\% & 21.1\% & 11.7\% & 12.1\% & 34.6\% & 16.5\% & 32.3\% \\
        \resultgroup{11}{空间推理模型}
        SpaceR & \cellcolor{secondblue}{28.0\%} & 24.0\% & 0.0\% & 2.6\% & 26.3\% & 15.3\% & 13.9\% & 33.7\% & 15.7\% & 30.8\% \\
        SpaceThinker & 25.0\% & 19.0\% & 0.0\% & 1.6\% & 22.4\% & 18.1\% & 14.1\% & 33.7\% & 13.6\% & 27.1\% \\
        SpaceOm & \cellcolor{bestblue}{30.9\%} & 33.2\% & 0.0\% & 3.2\% & 27.0\% & 16.7\% & 14.0\% & 32.3\% & 17.0\% & 34.2\% \\
        ViLaSR & 17.0\% & 5.3\% & 0.0\% & 0.8\% & 9.7\% & 3.4\% & 5.6\% & 35.9\% & 10.1\% & 20.6\% \\
        VST-7B-RL & 7.6\% & 0.0\% & 0.0\% & 0.0\% & 0.0\% & 0.0\% & 0.0\% & 22.8\% & 0.0\% & 0.0\% \\
        VST-7B-SFT & 7.8\% & 0.0\% & 0.0\% & 0.0\% & 0.0\% & 0.0\% & 0.0\% & 23.4\% & 4.1\% & 8.1\% \\
        \bottomrule
      \end{tabular}%
    }
\end{table}

从\autoref{tab:video-results}可以看出，视频榜单中 Gemini 3 Flash、Mimo v2 Omni、InternVL 3.5-38B、Qwen 3.5-Flash 等模型表现更有优势，但优势来源并不一致。Gemini 3 Flash 的视频总体准确率为 46.5\%，复合行为语义、原子行为语义和可见次数准确率都处于较高水平，说明其在视频语义和时间证据之间较为均衡。Mimo v2 Omni 的视频总体准确率为 38.8\%，尤其在可见次数准确率 59.4\% 上表现突出，说明它更擅长利用连续时间线索完成目标可见性判断。InternVL 3.5-38B 的可见次数准确率达到 64.0\%，但复合行为准确率只有 15.2\%，体现出识别显著目标出现次数与理解复杂飞行行为之间的能力分化。Qwen 3.5-4B 虽然视频总体准确率只有 16.5\%，但完整细粒度结果显示其可见区间指标仍存在局部能力，说明低综合分模型也可能在局部证据指标上有可挖掘空间。

视频分支的核心问题与图像分支平行，但证据形式由空间定位变为时间定位。选取各模型家族的代表性模型进行统计，视频任务中的平均语义成绩达到 33.20\%，时间证据指标 Temporal F1@0.5 的平均值为 29.02\%。从均值上看，视频分支中的语义与证据差距小于图像分支，但这并不意味着问题已经解决。恰恰相反，视频分支暴露出更复杂的“证据到语义映射不稳定”问题。一部分模型的时间定位指标甚至高于语义识别指标，这意味着模型可以大致察觉关键片段发生在何处，却未必能把这些片段正确归纳为复合行为语义或稳定的可见性语义；另一部分模型则表现为语义成绩高于时间定位成绩，说明它们能够给出较粗粒度的高层判断，但尚不能精确圈定支撑该判断的时间区间。

\begin{figure}[htbp]
    \centering
    \includegraphics[width=0.85\textwidth]{fig_analysis_video_semantic_vs_temporal.png}
    \caption{视频任务语义成绩与时间定位指标对比}
    \label{fig:video-semantic-temporal}
\end{figure}

从任务级相关性看，视频分支内部也呈现明显差异。\autoref{tab:video-correlation}显示，地标可见次数与区间推理任务的语义-时间相关性最高，Pearson \(r=0.84\)；原子行为识别次之，Pearson \(r=0.68\)；复合行为识别最低，Pearson \(r=0.53\)。

\begin{table}[htbp]
    \centering
    \caption{视频分支任务级相关性分析表}
    \label{tab:video-correlation}
    \begin{tabularx}{\textwidth}{Xrrr}
        \toprule
        任务 & Pearson \(r\) & 语义均值 & 时间证据均值 \\
        \midrule
        原子行为识别与时间定位 & 0.68 & 27.19 & 34.09 \\
        复合行为识别与时间定位 & 0.53 & 23.13 & 15.95 \\
        地标可见次数与区间推理 & 0.84 & 44.91 & 31.29 \\
        \bottomrule
    \end{tabularx}
\end{table}

该结果说明，视频分支的时间证据恢复整体优于图像分支的空间定位，但复合行为识别仍然表现出较明显的语义-区间脱节问题。也就是说，模型更容易识别目标“出现了几次”或某些原子动作“发生在何处”，却更难把多个局部片段整合为完整的高层飞行行为。这说明长时序行为语义整合仍是当前多模态大模型的突出短板。

\FloatBarrier
\subsubsection{跨媒介能力迁移分析}

为了进一步比较模型在两种媒介下的整体稳定性，\autoref{tab:combined-results}汇总了图像、视频、自我状态认知和环境情况认知四个聚合维度。该表既是综合榜单，也为跨媒介迁移分析和双认知差异分析提供共同参照。

\begin{table}[htbp]
    \centering
    \caption{综合实验结果}
    \label{tab:combined-results}
    \scriptsize
    \renewcommand{\arraystretch}{0.9}
    \setlength{\tabcolsep}{4pt}
    \resizebox{\textwidth}{!}{%
    \begin{tabular}{lrrrrr}
        \toprule
        模型 & 全局准确率 & \multicolumn{2}{c}{媒介分类} & \multicolumn{2}{c}{能力分类} \\
         &  & 图像 & 视频 & 自我状态认知 & 环境情况认知 \\
        \midrule
        \resultgroup{6}{闭源模型}
        Gemini 3 Flash & \cellcolor{bestblue}{48.4\%} & 50.2\% & 46.5\% & 44.4\% & 54.2\% \\
        Gemini 3.1 Flash Lite & 33.2\% & 34.5\% & 31.9\% & 27.5\% & 41.3\% \\
        Qwen 3.6-Plus & 34.5\% & 39.3\% & 29.6\% & 26.1\% & 47.2\% \\
        Qwen 3.5-Plus & 33.5\% & 38.5\% & 28.6\% & 22.5\% & 49.9\% \\
        Qwen 3.5-Flash & \cellcolor{secondblue}{38.9\%} & 40.9\% & 36.8\% & 28.7\% & 53.1\% \\
        Mimo v2 Omni & 35.0\% & 31.2\% & 38.8\% & 26.7\% & 44.7\% \\
        \resultgroup{6}{开源模型}
        GLM 4.6V & 32.6\% & 32.8\% & 32.4\% & 30.9\% & 34.9\% \\
        Kimi K2.5 & 32.5\% & 34.1\% & 30.9\% & 27.3\% & 40.0\% \\
        Qwen 3.5-397B-A17B & 31.4\% & 39.4\% & 23.4\% & 23.3\% & 44.9\% \\
        Qwen 3.5-122B-A10B & 36.4\% & 42.3\% & 30.6\% & 28.4\% & 49.1\% \\
        Qwen 3.5-35B-A3B & \cellcolor{bestblue}{37.9\%} & 41.9\% & 34.0\% & 26.2\% & 54.9\% \\
        Qwen 3.5-27B & \cellcolor{secondblue}{37.8\%} & 44.3\% & 31.3\% & 27.8\% & 53.3\% \\
        Qwen 3.5-9B & 35.1\% & 40.9\% & 29.2\% & 23.8\% & 52.1\% \\
        Qwen 3.5-4B & 27.8\% & 39.0\% & 16.5\% & 19.2\% & 42.9\% \\
        InternVL 3.5-30B-A3B & 30.0\% & 32.0\% & 28.0\% & 22.9\% & 40.2\% \\
        InternVL 3.5-14B & 33.9\% & 31.1\% & 36.6\% & 27.5\% & 41.4\% \\
        InternVL 3.5-8B & 30.6\% & 27.1\% & 34.1\% & 25.9\% & 35.8\% \\
        InternVL 3.5-4B & 28.3\% & 27.8\% & 28.9\% & 26.0\% & 31.2\% \\
        \resultgroup{6}{空间推理模型}
        SpaceR & \cellcolor{bestblue}{31.0\%} & 34.0\% & 28.0\% & 23.7\% & 41.6\% \\
        SpaceThinker & 26.9\% & 28.8\% & 25.0\% & 22.3\% & 33.6\% \\
        SpaceOm & \cellcolor{secondblue}{30.1\%} & 29.3\% & 30.9\% & 27.4\% & 33.4\% \\
        ViLaSR & 25.2\% & 33.4\% & 17.0\% & 14.0\% & 42.9\% \\
        VST-7B-RL & 21.9\% & 36.1\% & 7.6\% & 13.3\% & 38.0\% \\
        VST-7B-SFT & 21.9\% & 36.1\% & 7.8\% & 12.8\% & 38.8\% \\
        \bottomrule
      \end{tabular}%
    }
\end{table}

从\autoref{tab:combined-results}可以看出，综合榜单中 Gemini 3 Flash、Qwen 3.5-Flash、Qwen 3.5-35B-A3B 和 Qwen 3.5-27B 处于领先位置，说明当前较强模型已经能够在图像和视频两种媒介下同时获得较高的双认知总体准确率。但该表同时揭示，综合分相近的模型可能具有完全不同的媒介结构。例如，Mimo v2 Omni 的图像总体准确率为 31.18、视频总体准确率为 38.76，更偏向视频；Qwen 3.5-4B 的图像总体准确率为 38.99\%、视频总体准确率仅为 16.54\%，媒介差距明显；InternVL 3.5-14B 和 InternVL 3.5-8B 都是视频总体准确率高于图像总体准确率，说明 InternVL 系列在视频侧更具优势。由此可见，综合分不能替代媒介维度分析。

\begin{table}[htbp]
    \centering
    \caption{跨媒介相关性分析表}
    \label{tab:media-correlation}
    \scriptsize
    \setlength{\tabcolsep}{3pt}
    \begin{tabularx}{\textwidth}{@{}>{\raggedright\arraybackslash}X*{7}{>{\centering\arraybackslash}p{1.2cm}}@{}}
        \toprule
        比较项 & 配对模型数 & Pearson \(r\) & 图像均值 & 视频均值 & 均值差（视频-图像） & 视频更高模型数 & 图像更高模型数 \\
        \midrule
        图像总体准确率 vs 视频总体准确率 & 23.00 & 0.59 & 19.48 & 23.09 & 3.62 & 18.00 & 5.00 \\
        \bottomrule
    \end{tabularx}
\end{table}

\autoref{tab:media-correlation}显示，图像总体准确率与视频总体准确率之间的 Pearson 相关系数为 0.59，属于中等程度正相关；在 23 个配对模型中，有 18 个模型的视频总体准确率高于图像总体准确率，仅 5 个模型图像总体准确率更高。该结果说明，当前模型的跨媒介迁移是真实存在的：图像侧较强的模型往往在视频侧也不会太弱，因而图像与视频之间并非完全断裂的两套能力空间。

\begin{figure}[htbp]
    \centering
    \includegraphics[width=0.85\textwidth]{fig_analysis_image_vs_video_scatter.png}
    \caption{图像总体准确率与视频总体准确率的跨媒介散点图}
    \label{fig:image-video-scatter}
\end{figure}

结合\autoref{fig:image-video-scatter}可以看到，这种迁移仍远未达到稳定一致的程度。散点没有沿对角线紧密分布，而是表现出明显离散和偏移；视频均值 23.09\% 高于图像均值 19.48\%，均值差达到 3.62 个百分点，说明当前模型更容易在时间连续的证据环境下形成某种粗粒度判断。但这种视频侧优势并不意味着模型已经真正解决双认知问题，而更可能说明连续动态线索为模型提供了判断捷径。真正可靠的无人机智能体仍需要在单帧空间定位、跨视角匹配、长时序行为理解和时间区间证据之间保持一致。

定性案例同样支持这一判断。\autoref{fig:qualitative-case}展示了两类代表性失败模式：左侧为自我状态认知图像任务，模型需要从地标中心参考系转换到当前无人机视角。GPT 5.3 Chat 和 Claude Sonnet 4.6 大致定位了目标地标的位置，但没有从地标相关参考系出发完成方向转换，因此方向选项错误；Gemini 3.1 Flash Lite 同样没有完成参考坐标系转换，同时未输出有效的定位框；Qwen 3.5-Plus 判断正确方向，但目标框仍不够准确。右侧为环境情况认知视频任务，Qwen 3.5-Plus 判断正确可见次数，但时间区间不完整；GPT 5.3 Plus 和 Gemini 3.1 Flash Lite 未能连续定位到目标地标，导致次数判断错误；Claude Sonnet 4.6 能够持续、完整定位地标，表现最好。该案例说明，模型错误并不只是“看不见目标”，而是经常发生在语义判断、参考系转换和证据定位之间的连接环节。

\begin{figure}[htbp]
    \centering
    \includegraphics[width=\textwidth]{fig_qualitative_analysis.pdf}
    \caption{定性分析示例：左侧为自我状态认知图像任务失败案例，右侧为环境情况认知视频任务失败案例}
    \label{fig:qualitative-case}
\end{figure}

这些案例表明，当前模型失败并不只是简单答错，而更多表现为三类结构性不足：其一，模型无法稳定处理主视图、当前视图和未来视图之间的跨视角映射；其二，高层语义判断无法自然落到可检验的目标框或时间区间上；其三，模型能够察觉局部时间线索，却不能稳定整合为复合行为或完整事件过程。

\FloatBarrier
\subsection{双认知差异分析：当前模型明显偏向环境情况认知而非自我状态认知}

如果说跨媒介分析回答的是同一能力能否在图像与视频之间保持稳定，那么双认知差异分析回答的是模型是否同样理解无人机自身和外部环境。UAV-DualCog 的核心思想并不是把无人机任务简单看作视觉问答，而是将其拆分为自我状态认知和环境情况认知两条能力主线。

\subsubsection{自我状态认知能力结果与分析}

自我状态认知任务包括图像侧的地标相对位置推理、未来观察预测，以及视频侧的复合行为识别与原子行为识别。它们共同要求模型理解无人机自己在哪里、会看到什么、正在怎样运动，比单纯识别外部目标更依赖主体状态建模、参考坐标系转换和视角变化理解。\autoref{tab:self-cognition-results}给出了自我状态认知排行榜。

\begin{table}[htbp]
    \centering
    \caption{自我状态认知排行榜完整结果}
    \label{tab:self-cognition-results}
    \scriptsize
    \renewcommand{\arraystretch}{0.9}
    \setlength{\tabcolsep}{4pt}
    \resizebox{\textwidth}{!}{
    \begin{tabular}{lrrrrrrrrrrrrr}
        \toprule
        模型 & 全局准确率 & \multicolumn{6}{c}{图像任务} & \multicolumn{6}{c}{视频任务} \\
         &  & \multicolumn{3}{c}{地标相对位置推理} & \multicolumn{3}{c}{未来观察预测} & \multicolumn{3}{c}{复合行为识别} & \multicolumn{3}{c}{原子行为识别} \\
         &  & Acc & BBox Acc@0.5 & mIoU & Acc & BBox Acc@0.5 & mIoU & Acc & F1@0.5 & Mean tIoU & Acc & F1@0.5 & Mean tIoU \\
        \midrule
        \resultgroup{14}{闭源模型}
        Gemini 3 Flash & \cellcolor{bestblue}{44.4\%} & 47.6\% & 0.7\% & 1.2\% & 45.9\% & 0.7\% & 0.9\% & 49.3\% & 51.9\% & 54.3\% & 35.0\% & 52.6\% & 46.6\% \\
        Gemini 3.1 Flash Lite & 27.5\% & 39.1\% & 0.0\% & 0.0\% & 28.6\% & 2.5\% & 3.5\% & 22.3\% & 29.5\% & 32.2\% & 19.7\% & 40.0\% & 40.6\% \\
        Qwen 3.6-Plus & 26.1\% & 32.8\% & 16.6\% & 27.1\% & 33.5\% & 14.0\% & 26.7\% & 9.8\% & 16.5\% & 16.9\% & 28.3\% & 45.9\% & 47.3\% \\
        Qwen 3.5-Plus & 22.5\% & 28.6\% & 21.4\% & 27.1\% & 29.0\% & 13.4\% & 23.5\% & 2.4\% & 9.0\% & 9.2\% & 30.0\% & 46.6\% & 47.7\% \\
        Qwen 3.5-Flash & \cellcolor{secondblue}{28.7\%} & 27.8\% & 21.6\% & 29.8\% & 28.9\% & 15.9\% & 23.3\% & 32.6\% & 35.8\% & 40.1\% & 25.6\% & 51.7\% & 42.4\% \\
        Mimo v2 Omni & 26.7\% & 28.7\% & 38.0\% & 34.3\% & 21.3\% & 29.8\% & 29.7\% & 27.8\% & 33.1\% & 35.8\% & 29.1\% & 53.6\% & 46.9\% \\
        \resultgroup{14}{开源模型}
        GLM 4.6V & \cellcolor{bestblue}{30.9\%} & 30.8\% & 9.4\% & 8.7\% & 27.1\% & 8.6\% & 8.7\% & 32.4\% & 1.3\% & 6.7\% & 33.4\% & 21.6\% & 22.6\% \\
        Kimi K2.5 & 27.3\% & 34.3\% & 27.1\% & 30.4\% & 32.7\% & 15.0\% & 21.7\% & 11.9\% & 19.1\% & 19.9\% & 30.2\% & 54.6\% & 49.2\% \\
        Qwen 3.5-397B-A17B & 23.3\% & 27.4\% & 3.3\% & 12.6\% & 35.8\% & 10.0\% & 21.4\% & 5.2\% & 1.7\% & 2.2\% & 24.6\% & 29.9\% & 30.0\% \\
        Qwen 3.5-122B-A10B & \cellcolor{secondblue}{28.4\%} & 32.6\% & 29.1\% & 31.6\% & 37.8\% & 20.6\% & 28.4\% & 21.2\% & 28.4\% & 31.9\% & 22.1\% & 39.0\% & 34.6\% \\
        Qwen 3.5-35B-A3B & 26.2\% & 29.5\% & 27.6\% & 30.5\% & 27.4\% & 0.5\% & 0.8\% & 21.6\% & 21.5\% & 23.1\% & 26.2\% & 54.4\% & 50.3\% \\
        Qwen 3.5-27B & 27.8\% & 31.6\% & 42.8\% & 39.9\% & 37.1\% & 17.1\% & 25.8\% & 14.6\% & 21.9\% & 23.3\% & 28.0\% & 54.5\% & 50.7\% \\
        Qwen 3.5-9B & 23.8\% & 29.0\% & 15.9\% & 25.8\% & 29.0\% & 17.8\% & 29.8\% & 9.9\% & 10.1\% & 11.0\% & 27.2\% & 55.9\% & 51.3\% \\
        Qwen 3.5-4B & 19.2\% & 30.9\% & 12.8\% & 23.1\% & 30.2\% & 11.2\% & 21.7\% & 6.4\% & 1.3\% & 1.4\% & 9.4\% & 3.0\% & 2.8\% \\
        InternVL 3.5-30B-A3B & 22.9\% & 27.3\% & 4.9\% & 19.9\% & 24.0\% & 7.0\% & 15.5\% & 15.3\% & 5.8\% & 8.6\% & 24.8\% & 24.2\% & 24.6\% \\
        InternVL 3.5-14B & 27.5\% & 22.8\% & 7.6\% & 22.2\% & 31.9\% & 6.6\% & 22.0\% & 33.1\% & 3.9\% & 11.3\% & 22.4\% & 15.8\% & 17.4\% \\
        InternVL 3.5-8B & 25.9\% & 28.5\% & 19.7\% & 34.2\% & 24.5\% & 8.1\% & 28.7\% & 27.4\% & 8.5\% & 14.1\% & 23.0\% & 25.0\% & 26.1\% \\
        InternVL 3.5-4B & 26.0\% & 25.6\% & 3.1\% & 19.8\% & 26.4\% & 10.8\% & 23.2\% & 31.0\% & 3.5\% & 7.4\% & 21.1\% & 11.7\% & 12.1\% \\
        \resultgroup{14}{空间推理模型}
        SpaceR & \cellcolor{secondblue}{23.7\%} & 20.6\% & 2.4\% & 6.0\% & 24.0\% & 1.2\% & 1.9\% & 24.0\% & 0.0\% & 2.6\% & 26.3\% & 15.3\% & 13.9\% \\
        SpaceThinker & 22.3\% & 22.6\% & 1.2\% & 3.0\% & 25.4\% & 1.9\% & 5.7\% & 19.0\% & 0.0\% & 1.6\% & 22.4\% & 18.1\% & 14.1\% \\
        SpaceOm & \cellcolor{bestblue}{27.4\%} & 24.7\% & 2.8\% & 6.7\% & 24.5\% & 3.6\% & 10.9\% & 33.2\% & 0.0\% & 3.2\% & 27.0\% & 16.7\% & 14.0\% \\
        ViLaSR & 14.0\% & 18.8\% & 0.0\% & 0.0\% & 22.1\% & 0.0\% & 0.0\% & 5.3\% & 0.0\% & 0.8\% & 9.7\% & 3.4\% & 5.6\% \\
        VST-7B-RL & 13.3\% & 29.2\% & 0.0\% & 0.0\% & 24.1\% & 0.0\% & 0.0\% & 0.0\% & 0.0\% & 0.0\% & 0.0\% & 0.0\% & 0.0\% \\
        VST-7B-SFT & 12.8\% & 28.1\% & 0.0\% & 0.0\% & 23.0\% & 0.0\% & 0.0\% & 0.0\% & 0.0\% & 0.0\% & 0.0\% & 0.0\% & 0.0\% \\
        \bottomrule
      \end{tabular}}
\end{table}

从\autoref{tab:self-cognition-results}可以看出，自我状态认知榜单整体分数偏低。Gemini 3 Flash 以 44.44 排名第一，但这已经显著低于它的环境情况认知全局准确率 54.20；排名第二的 GLM 4.6V 只有 30.91，其后大量模型集中在 20 多的区间。这表明当前模型对自身状态的建模能力明显不足。进一步看，不同模型在自我状态认知内部也呈现出明显分化。Qwen 3.5-27B 在地标相对位置推理 BBox Acc@0.5 上达到 42.77，说明它能够在部分空间自我定位任务中给出较好证据；但其复合行为语义只有 14.62，说明空间自我定位能力并不能自然迁移到复杂行为理解。Qwen 3.5-9B 的原子行为 F1@0.5 达到 55.92，但复合行为语义只有 9.87，说明模型可能捕捉到局部运动片段，却难以把局部动作组合成高层飞行行为。由此可见，自我状态认知不是单一能力，而是由跨视角空间映射、运动后观察预测、行为识别和时间定位共同构成的复合能力。

\subsubsection{环境情况认知能力结果与分析}

环境情况认知任务包括图像侧的自我参照位置推理、地标驱动动作决策，以及视频侧的地标可见次数与区间推理。与自我状态认知相比，环境情况认知更关注外界目标在哪里、应如何根据目标位置行动、目标在视频中何时可见。\autoref{tab:environment-cognition-results}给出了环境情况认知排行榜。

\begin{table}[htbp]
    \centering
    \caption{环境情况认知排行榜完整结果}
    \label{tab:environment-cognition-results}
    \scriptsize
    \renewcommand{\arraystretch}{0.9}
    \setlength{\tabcolsep}{4pt}
    \resizebox{\textwidth}{!}{
    \begin{tabular}{lrrrrrrrrrr}
        \toprule
        模型 & 全局准确率 & \multicolumn{6}{c}{图像任务} & \multicolumn{3}{c}{视频任务} \\
         &  & \multicolumn{3}{c}{自我参照位置推理} & \multicolumn{3}{c}{地标驱动动作决策} & \multicolumn{3}{c}{地标可见性推理} \\
         &  & Acc & BBox Acc@0.5 & mIoU & Acc & BBox Acc@0.5 & mIoU & Acc & F1@0.5 & Mean tIoU \\
        \midrule
        \resultgroup{11}{闭源模型}
        Gemini 3 Flash & \cellcolor{bestblue}{54.2\%} & 56.2\% & 0.1\% & 0.9\% & 51.3\% & 0.5\% & 1.2\% & 55.2\% & 44.9\% & 44.3\% \\
        Gemini 3.1 Flash Lite & 41.3\% & 39.9\% & 0.4\% & 0.8\% & 30.4\% & 0.4\% & 0.6\% & 53.7\% & 43.1\% & 41.6\% \\
        Qwen 3.6-Plus & 47.2\% & 48.4\% & 7.4\% & 19.0\% & 42.5\% & 6.0\% & 19.9\% & 50.8\% & 51.5\% & 46.1\% \\
        Qwen 3.5-Plus & \cellcolor{secondblue}{49.9\%} & 47.6\% & 9.7\% & 19.3\% & 48.9\% & 8.3\% & 20.0\% & 53.2\% & 53.9\% & 46.8\% \\
        Qwen 3.5-Flash & 53.1\% & 52.4\% & 8.8\% & 19.4\% & 54.5\% & 8.9\% & 20.1\% & 52.3\% & 57.8\% & 51.3\% \\
        Mimo v2 Omni & 44.7\% & 36.5\% & 8.8\% & 17.4\% & 38.2\% & 9.2\% & 17.3\% & 59.4\% & 52.5\% & 53.7\% \\
        \resultgroup{11}{开源模型}
        GLM 4.6V & 34.9\% & 29.2\% & 3.3\% & 6.7\% & 44.2\% & 2.5\% & 5.0\% & 31.3\% & 16.7\% & 33.8\% \\
        Kimi K2.5 & 40.0\% & 39.4\% & 5.0\% & 15.6\% & 30.1\% & 4.7\% & 14.4\% & 50.5\% & 49.7\% & 45.9\% \\
        Qwen 3.5-397B-A17B & 44.9\% & 47.9\% & 2.8\% & 13.0\% & 46.2\% & 4.4\% & 13.9\% & 40.4\% & 20.4\% & 31.8\% \\
        Qwen 3.5-122B-A10B & 49.1\% & 49.7\% & 6.1\% & 19.4\% & 49.0\% & 6.8\% & 18.8\% & 48.4\% & 46.4\% & 49.7\% \\
        Qwen 3.5-35B-A3B & \cellcolor{bestblue}{54.9\%} & 53.2\% & 6.1\% & 14.5\% & 57.5\% & 7.2\% & 18.5\% & 54.1\% & 49.7\% & 48.6\% \\
        Qwen 3.5-27B & \cellcolor{secondblue}{53.3\%} & 57.8\% & 7.3\% & 20.4\% & 50.6\% & 7.5\% & 21.0\% & 51.4\% & 51.7\% & 48.7\% \\
        Qwen 3.5-9B & 52.1\% & 53.2\% & 10.3\% & 20.9\% & 52.2\% & 11.5\% & 22.1\% & 50.7\% & 33.5\% & 36.8\% \\
        Qwen 3.5-4B & 42.9\% & 47.5\% & 6.7\% & 17.7\% & 47.5\% & 6.3\% & 18.0\% & 33.9\% & 44.3\% & 52.0\% \\
        InternVL 3.5-30B-A3B & 40.2\% & 35.4\% & 0.6\% & 7.6\% & 41.3\% & 1.3\% & 7.3\% & 43.9\% & 7.9\% & 12.6\% \\
        InternVL 3.5-14B & 41.4\% & 28.5\% & 4.7\% & 12.9\% & 41.3\% & 5.0\% & 11.9\% & 54.5\% & 13.9\% & 22.7\% \\
        InternVL 3.5-8B & 35.8\% & 24.7\% & 4.2\% & 12.3\% & 30.8\% & 2.5\% & 11.4\% & 51.9\% & 16.4\% & 21.3\% \\
        InternVL 3.5-4B & 31.2\% & 22.4\% & 1.0\% & 8.0\% & 36.7\% & 1.0\% & 7.5\% & 34.6\% & 16.5\% & 32.3\% \\
        \resultgroup{11}{空间推理模型}
        SpaceR & \cellcolor{secondblue}{41.6\%} & 42.4\% & 1.8\% & 5.4\% & 48.8\% & 2.3\% & 6.8\% & 33.7\% & 15.7\% & 30.8\% \\
        SpaceThinker & 33.6\% & 32.8\% & 0.6\% & 5.1\% & 34.3\% & 0.7\% & 4.7\% & 33.7\% & 13.6\% & 27.1\% \\
        SpaceOm & 33.4\% & 35.2\% & 2.4\% & 8.6\% & 32.8\% & 1.8\% & 6.5\% & 32.3\% & 17.0\% & 34.2\% \\
        ViLaSR & \cellcolor{bestblue}{42.9\%} & 43.1\% & 0.0\% & 0.0\% & 49.7\% & 0.0\% & 0.0\% & 35.9\% & 10.1\% & 20.6\% \\
        VST-7B-RL & 38.0\% & 47.2\% & 0.0\% & 0.1\% & 44.0\% & 0.0\% & 0.0\% & 22.8\% & 0.0\% & 0.0\% \\
        VST-7B-SFT & 38.8\% & 46.5\% & 0.0\% & 0.2\% & 46.6\% & 0.0\% & 0.2\% & 23.4\% & 4.1\% & 8.1\% \\
        \bottomrule
      \end{tabular}}
\end{table}

从\autoref{tab:environment-cognition-results}可以看出，环境情况认知榜单整体显著高于自我状态认知榜单。排名前五的模型环境情况认知全局总体准确率均超过 52，其中 Qwen 3.5-35B-A3B 达到 54.95，Gemini 3 Flash 达到 54.20，Qwen 3.5-27B 和 Qwen 3.5-Flash 也分别达到 53.26 和 53.09。这说明当前模型更容易处理外部目标的位置、动作决策和可见性过程。

但环境情况认知的高分同样需要细分理解。Gemini 3 Flash 的环境情况认知全局总体准确率很高，但图像侧自我参照位置推理 BBox Acc@0.5 只有 0.10、地标驱动动作决策 BBox Acc@0.5 只有 0.49，说明它的环境情况认知优势主要来自选项判断和视频可见性，而不是精细空间证据。Qwen 3.5-Flash 的环境情况认知更均衡，自我参照、动作决策和可见性区间指标都较高，其中可见区间 F1@0.5 达到 57.84。Mimo v2 Omni 的图像侧环境任务并不突出，但可见次数准确率 59.38、可见区间 F1@0.5 52.52 表现很强，说明它主要依赖视频连续线索形成环境情况认知优势。由此可见，环境情况认知高于自我状态认知并不意味着模型已经完全理解外界，而是说明外界目标识别、可见性计数和粗粒度方位判断相对更容易被当前模型掌握。

\FloatBarrier
\subsubsection{双认知能力差异分析}

为了更直观地比较同一模型在两条能力主线上的分布差异，本文先给出自我状态认知与环境情况认知的综合结果，再进一步统计二者之间的相关性与均值差距。

\begin{table}[!htbp]
    \centering
    \caption{自我状态认知与环境情况认知综合结果}
    \label{tab:cognition-results}
    \small
    \renewcommand{\arraystretch}{0.95}
    \setlength{\tabcolsep}{5pt}
    \begin{tabularx}{\textwidth}{@{}>{\raggedright\arraybackslash}X>{\centering\arraybackslash}p{2.3cm}>{\centering\arraybackslash}p{2.8cm}>{\centering\arraybackslash}p{2.8cm}@{}}
        \toprule
        模型 & 全局准确率 & 自我状态认知 & 环境情况认知 \\
        \midrule
        Gemini 3 Flash & 48.4\% & 44.4\% & 54.2\% \\
        Qwen 3.5-Flash & 38.9\% & 28.7\% & 53.1\% \\
        Qwen 3.5-35B-A3B & 37.9\% & 26.2\% & 54.9\% \\
        Qwen 3.5-27B & 37.8\% & 27.8\% & 53.3\% \\
        Qwen 3.5-122B-A10B & 36.4\% & 28.4\% & 49.1\% \\
        Mimo v2 Omni & 35.0\% & 26.7\% & 44.7\% \\
        Kimi K2.5 & 32.5\% & 27.3\% & 40.0\% \\
        GLM 4.6V & 32.6\% & 30.9\% & 34.9\% \\
        SpaceOm & 30.1\% & 27.4\% & 33.4\% \\
        ViLaSR & 25.2\% & 14.0\% & 42.9\% \\
        VST-7B-RL & 21.9\% & 13.3\% & 38.0\% \\
        \bottomrule
    \end{tabularx}
\end{table}

结合\autoref{tab:combined-results}和\autoref{tab:cognition-results}可以看到，综合榜前列模型同样呈现结构性失衡：Gemini 3 Flash 的自我状态认知为 44.44、环境情况认知为 54.20，Qwen 3.5-Flash 分别为 28.73 和 53.09，Qwen 3.5-35B-A3B 分别为 26.19 和 54.95。也就是说，较高综合分并不意味着模型具备均衡的双认知能力；很多模型之所以排名靠前，很大程度上来自环境情况认知分数的支撑。若没有双认知框架，模型在“理解自身”这一核心能力上的短板很容易被总体分数掩盖。

\begin{table}[!htbp]
    \centering
    \caption{双认知平衡相关性分析表}
    \label{tab:cognition-correlation}
    \scriptsize
    \setlength{\tabcolsep}{4pt}
    \begin{tabularx}{\textwidth}{@{}>{\raggedright\arraybackslash}X>{\centering\arraybackslash}p{1.4cm}>{\centering\arraybackslash}p{1.4cm}>{\centering\arraybackslash}p{1.7cm}>{\centering\arraybackslash}p{1.9cm}>{\centering\arraybackslash}p{1.9cm}>{\centering\arraybackslash}p{1.7cm}>{\centering\arraybackslash}p{1.7cm}@{}}
        \toprule
        比较项 & 配对模型数 & Pearson \(r\) & 自我状态认知均值 & 环境情况认知均值 & 均值差（环境-自我） & 环境更高模型数 & 自我更高模型数 \\
        \midrule
        自我状态认知总体准确率 vs 环境情况认知总体准确率 & 24.00 & 0.34 & 26.26 & 43.24 & 16.97 & 24.00 & 0.00 \\
        \bottomrule
    \end{tabularx}
\end{table}

\autoref{tab:cognition-correlation}显示，自我状态认知总体准确率与环境情况认知总体准确率之间的 Pearson 相关系数仅为 0.34，说明两类能力之间只有有限正相关；同时，自我状态认知均值为 26.26\%，环境情况认知均值达到 43.24\%，二者相差 16.97 个百分点。在当前 24 个配对模型中，全部模型都是环境情况认知高于自我状态认知。

\begin{figure}[!htbp]
    \centering
    \includegraphics[width=0.85\textwidth]{fig_analysis_self_vs_environment_scatter.png}
    \caption{自我状态认知与环境情况认知总体准确率散点图}
    \label{fig:self-env-scatter}
\end{figure}

\autoref{fig:self-env-scatter}进一步说明，当前模型更容易理解外部目标与环境，却更难真正理解无人机自身。环境情况认知任务更多依赖目标外观、方位关系和可见性判断；自我状态认知任务则要求模型同时处理参考坐标系、无人机视角变化、飞行动作模式以及行为阶段，更接近具身智能中的主体建模问题。自我状态认知由此成为当前模型最明显的短板。

\FloatBarrier
\subsection{深入交叉分析}

综合跨媒介分析和双认知差异分析，可以得到进一步发现。

首先，不同模型家族呈现明显不同的优势和短板，综合得分不能替代能力结构分析。闭源模型中，Gemini 3 Flash 以 48.35 的综合分位居首位，并在图像 50.22、视频 46.49 两个媒介上都保持较强表现，但其图像任务中的多个 BBox Acc@0.5 指标接近 0，说明强语义模型仍可能缺乏精细空间证据。GPT 5.3 Chat、Claude Sonnet 4.6 在图像选项准确率上较强，但闭源模型整体也没有自然解决证据对齐问题。开源模型中，Qwen 3.5 系列表现最稳定，Qwen 3.5-Flash、Qwen 3.5-35B-A3B、Qwen 3.5-27B 的环境情况认知总体准确率均超过 53，其中 Qwen 3.5-27B 在地标相对位置推理的 BBox Acc@0.5 达到 42.77，体现出较好的局部空间定位潜力。InternVL 系列则更偏向视频侧，InternVL 3.5-14B 和 InternVL 3.5-8B 的视频成绩均高于图像成绩，InternVL 3.5-38B 在视频榜中也具有较强竞争力。Mimo v2 Omni 在视频总体准确率和可见性区间指标上突出，说明部分模型更擅长利用连续时间线索。由此可见，模型家族差异不仅是“谁更强”，更是“强在语义、空间证据、时间证据还是媒介适应性”的差异。

其次，模型规模、MoE 架构和已有空间推理微调方法并不能自动转化为 UAV-DualCog 上的可靠双认知能力。以 Qwen 系列为例，MoE 模型并未随参数规模单调提升：Qwen 3.5-35B-A3B 综合分为 37.95，与 Dense 架构的 Qwen 3.5-27B 的 37.81 基本相当，而更大规模的 Qwen 3.5-122B-A10B 和 Qwen 3.5-397B-A17B 分别为 36.44 和 31.38。这说明路由容量和通用知识规模并不能直接弥合无人机双认知中的参照系、证据定位和时序整合问题。空间推理微调模型也呈现类似现象：ViLaSR、VST-7B-RL 等模型在部分选项任务上并非完全失效，但图像 BBox 指标大量为 0，说明一般空间推理或视觉语言后训练与 UAV-DualCog 所要求的“选项判断+显式证据”结构并不完全对齐。因此，后续优化不能简单依赖已有空间模型或更大模型，而需要围绕本课题定义的证据约束进行任务内训练。

再次，从零样本模型和微调模型对比来看，当前强通用模型在零样本设置下已经能够获得较高选项准确率，但空间或时间证据并不稳定；已有空间推理或视觉语言后训练模型虽然经过特定能力强化，却没有在 UAV-DualCog 中自动获得优势。这说明 UAV-DualCog 要求的不是一般意义上的空间推理或通用视觉问答能力，而是面向无人机场景、双认知任务和结构化证据输出共同对齐的任务内能力。

最后，能力评估结果也启发了后续能力优化的技术路线：训练目标应从提高答案准确率转向同时缩小语义-证据差距和自我-环境差距。当前实验表明，双认知能力不足主要体现在三方面：一是自我状态认知长期弱于环境情况认知，二是图像侧“语义相对较强、空间证据明显偏弱”的现象尤为突出，三是图像与视频之间尚未形成稳定一致的双认知能力表示。因此，优化不能只追求综合分数提升，而应围绕自我状态认知补强、语义-证据一致性强化、跨媒介稳定性提升三个目标展开。第4章已经据此完成对应训练方法、最终模型评测和结果分析。

UAV-DualCog 的任务输出是高度结构化的：图像任务要求输出“选项+目标框”，视频任务要求输出“语义结果+时间区间”。这使得监督微调适合作为基础训练方式，因为它能够通过统一的 next-token prediction 目标，让模型先学会稳定生成“答案+证据”的结构化输出。但当前实验也说明，只学习输出格式并不足以充分解决更细粒度的证据一致性问题，因为很多模型已经能够输出格式正确的结果，却没有真正依赖正确证据完成判断。因此，第4章进一步引入联合 GRPO，围绕答案正确性、格式合法性以及目标框质量进行有针对性的后训练优化，并通过最终模型评测验证其收益边界。

选择 Qwen 3.5 4B 作为优化基座模型具有明确问题导向：它的图像成绩为 38.99，说明具备一定空间问答基础，但视频成绩和自我状态认知仍有明显提升空间；同时 4B 规模适合在现有算力下开展多模态 SFT、LoRA 和 GRPO 实验。第4章的最终评测结果进一步表明，针对四类图像任务进行证据感知训练后，最终导出的 GRPO 权重模型能够显著提升图像侧选项准确率和空间证据质量，但这种收益在视频侧只对自我行为时间定位形成正向迁移，对环境可见性推理尚未形成稳定改善。

\subsection{能力评测结论}

当前实验可以形成明确结论：多模态大模型已经能够在无人机场景中捕捉部分目标、方位和动态线索，但尚未形成稳定、均衡、可证据验证的无人机双认知能力。真正面向具身无人机的多模态推理评测不能只问模型“语义答案是否正确”，还必须追问模型是否知道自己在何处、是否理解外界如何变化、是否能够给出支撑判断的空间或时间证据。

UAV-DualCog 的价值在于，它把多模态大模型的自我建模、环境理解、显式证据和跨媒介稳定性放入同一评测框架中。实验结果表明，当前模型明显偏向环境情况认知而非自我状态认知，图像和视频之间存在真实但不充分的迁移，语义正确性与证据可靠性之间也存在结构性差距。这些发现说明，面向无人机具身推理的可靠多模态系统需要实现自我状态认知与环境情况认知的匹配、协同发展，并在观测媒介变化时保持空间或时间证据的一致性。第4章在这一结论基础上给出能力优化方法，并进一步分析最终导出的 GRPO 权重模型在图像任务中的直接收益和在视频任务中的迁移分化。
